\documentclass[a4paper,12pt]{article}

% ======================
% Encodage / langue
% ======================
\usepackage[utf8]{inputenc}
\usepackage[T1]{fontenc}
\usepackage[french]{babel}

% ======================
% Mise en page / style
% ======================
\usepackage[margin=1in]{geometry}
\usepackage{graphicx}
\usepackage{xcolor}
\usepackage{tikz}
\usepackage{titlesec}
\usepackage{tocloft}
\usepackage{setspace}
\usepackage{float}
\usepackage{placeins}
\usepackage{xurl}
\usepackage{hyperref}
\usepackage{enumitem}
\usepackage{booktabs}

% ======================
% Couleurs
% ======================
\definecolor{ensiared}{RGB}{156,31,46}
\definecolor{tocblue}{RGB}{0,102,204}

% ======================
% ToC en bleu
% ======================
\renewcommand{\cftsecfont}{\color{tocblue}}
\renewcommand{\cftsubsecfont}{\color{tocblue}}
\renewcommand{\cftsecpagefont}{\color{tocblue}}
\renewcommand{\cftsubsecpagefont}{\color{tocblue}}

% ======================
% Interligne / spacing
% ======================
\linespread{1.15}
\titlespacing*{\section}{0pt}{*2}{*1.2}
\titlespacing*{\subsection}{0pt}{*1.5}{*1}

% ======================
% Encadré robuste
% ======================
\usepackage[most]{tcolorbox}
\newtcolorbox{tipbox}{
  colback=black!3,
  colframe=black!25,
  boxrule=0.6pt,
  arc=2pt,
  left=6pt,right=6pt,top=6pt,bottom=6pt
}

\pagestyle{empty}

\begin{document}

% ==========================================================
% Page de garde
% ==========================================================
\pagenumbering{gobble}

\begin{titlepage}
  \begin{tikzpicture}[remember picture, overlay]
    \fill[ensiared] ([yshift=2cm]current page.south west) rectangle (current page.south east);
    \fill[ensiared] ([yshift=2cm]current page.north west) rectangle (current page.north east);
  \end{tikzpicture}

  \centering
  \vspace*{2cm}

  % Logos
  \includegraphics[width=0.35\textwidth]{img/sd logo.jpg}\hfill
  \includegraphics[width=0.25\textwidth]{img/uca logo.png}\hfill

  \vspace*{1.0cm}

  {\color{white}\scshape\LARGE Base de données \& Systèmes décisionnels\par}
  \vspace{0.4cm}
  {\scshape\Large UNIVERSITÉ CLERMONT AUVERGNE — Campus Aurillac\par}

  \vspace{1.2cm}
  \hrule height 1pt
  \vspace{0.5cm}
  {\huge\bfseries Fiche technique — Outil décisionnel Ventes \& Logistique\par}
  \vspace{0.2cm}
  {\Large PostgreSQL + Datawarehouse + Power BI \;|\; SAE 6.VCOD.01\par}
  \vspace{0.5cm}
  \hrule height 1pt
  \vspace{1.6cm}

  \begin{minipage}[t]{0.5\textwidth}
    \raggedright
    {\Large\textbf{Réalisé par} :\par}
    \vspace{0.2cm}
    \begin{tabular}{l}
      Ibrahima BODIAN \\
    \end{tabular}
  \end{minipage}%
  \begin{minipage}[t]{0.5\textwidth}
    \raggedleft
    {\Large\textbf{Encadrant} :\par}
    \vspace{0.2cm}
    \textit{Bruno LAGARDE} \\
  \end{minipage}

  \vfill
  {\large Date : \textit{\today}\par}
\end{titlepage}

% ==========================================================
% Sommaire
% ==========================================================
\renewcommand{\contentsname}{Sommaire}
\tableofcontents
\newpage
\pagenumbering{arabic}
\setcounter{page}{1}

% ==========================================================
% Corps
% ==========================================================

\section*{1) Objectif technique}
\addcontentsline{toc}{section}{1) Objectif technique}
Mettre en place une chaîne décisionnelle complète :
\begin{itemize}
  \item import des données opérationnelles (OLTP) dans PostgreSQL (\texttt{ventebdd}, schéma \texttt{entreprise}) ;
  \item construction d’un datawarehouse en schéma en étoile (\texttt{entreprise\_dw}) ;
  \item alimentation du DW depuis l’OLTP via scripts SQL (\texttt{INSERT...SELECT}) ;
  \item restitution dans Power BI (Import, modèle en étoile, rapports Finance \textbf{F} et Logistique \textbf{L}).
\end{itemize}

\begin{tipbox}
\textbf{Remarque :} le guide utilisateur (1 page) décrit l’usage Power BI. Cette fiche décrit l’architecture, le modèle, le chargement et les contrôles.
\end{tipbox}

\section*{2) Environnement \& composants}
\addcontentsline{toc}{section}{2) Environnement \& composants}
\begin{itemize}
  \item \textbf{SGBD :} PostgreSQL (base \texttt{ventebdd})
  \item \textbf{Schémas :} \texttt{entreprise} (OLTP), \texttt{entreprise\_dw} (DW)
  \item \textbf{Sources :} fichiers CSV importés dans l’OLTP
  \item \textbf{BI :} Power BI Desktop (\textbf{Import})
\end{itemize}

\section*{3) Modèle source (OLTP) — schéma \texttt{entreprise}}
\addcontentsline{toc}{section}{3) Modèle source (OLTP) — schéma \texttt{entreprise}}
\textbf{Tables principales :} \texttt{\_client}, \texttt{\_commande}, \texttt{\_detailcommande}, \texttt{\_produit}, \texttt{\_categorie}, \texttt{\_fournisseur}, \texttt{\_employe}, \texttt{\_transporteur}.

\textbf{Relations clés :}
\begin{itemize}
  \item client $\rightarrow$ commande (1,n) ; commande $\rightarrow$ détailcommande (1,n) ;
  \item produit $\rightarrow$ détailcommande (1,n) ; catégorie $\rightarrow$ produit (1,n) ;
  \item fournisseur $\rightarrow$ produit (1,n) ; employé $\rightarrow$ commande (1,n) ;
  \item transporteur $\rightarrow$ commande (1,n).
\end{itemize}

\section*{4) Datawarehouse (DW) — schéma \texttt{entreprise\_dw}}
\addcontentsline{toc}{section}{4) Datawarehouse (DW) — schéma \texttt{entreprise\_dw}}
\textbf{Dimensions :}
\begin{itemize}
  \item \texttt{dimTemps} (date\_key, date\_jour), \texttt{dimClient}, \texttt{dimEmploye},
  \texttt{dimTransporteur}, \texttt{dimGeoLivraison}, \texttt{dimProduit} (dénormalisée : produit + catégorie + fournisseur).
\end{itemize}

\textbf{Faits :}
\begin{itemize}
  \item \texttt{faitVentes} (grain = ligne de commande) : qte, prixunit, remise, montants brut/net, \textit{port alloué}, nocom.
  \item \texttt{faitLivraisons} (grain = commande) : port, délais, écarts objectif/envoi, dates (com/envoi/objectif), nocom.
\end{itemize}

\section*{5) Règles de calcul (principes)}
\addcontentsline{toc}{section}{5) Règles de calcul (principes)}
\textbf{Ventes :}
\begin{itemize}
  \item Montant brut = \texttt{qte} $\times$ \texttt{prixunit}
  \item Montant net = \texttt{qte} $\times$ \texttt{prixunit} $\times (1-\texttt{remise})
\end{itemize}

\textbf{Port alloué :} le port (niveau commande) est réparti sur les lignes au prorata du montant net (évite le double comptage).

\textbf{Logistique :}
\begin{itemize}
  \item Délai d’envoi (jours) = \texttt{dateenv} $-$ \texttt{datecom}
  \item Écart objectif/envoi (jours) = \texttt{dateenv} $-$ \texttt{dateobjliv} (peut être négatif)
\end{itemize}

\section*{6) Alimentation DW (ETL SQL)}
\addcontentsline{toc}{section}{6) Alimentation DW (ETL SQL)}
\textbf{Ordre recommandé :}
\begin{enumerate}
  \item Dimensions : \texttt{dimTemps}, \texttt{dimClient}, \texttt{dimEmploye}, \texttt{dimTransporteur}, \texttt{dimGeoLivraison}, \texttt{dimProduit}
  \item Faits : \texttt{faitLivraisons} puis \texttt{faitVentes}
\end{enumerate}

\textbf{Nettoyage :} remplacement de \texttt{unknown\_value\_please\_contact\_support} par \textbf{NULL} lors du chargement des dimensions.

\section*{7) Contrôles qualité (cohérence)}
\addcontentsline{toc}{section}{7) Contrôles qualité (cohérence)}
\textbf{CA net — source vs DW :}
\begin{verbatim}
SELECT
  (SELECT SUM(qte * prixunit * (1 - remise))
   FROM entreprise._detailcommande) AS ca_oltp,
  (SELECT SUM(montant_net)
   FROM entreprise_dw.faitVentes) AS ca_dw;
\end{verbatim}

\textbf{Nb commandes — source vs DW :}
\begin{verbatim}
SELECT
  (SELECT COUNT(*) FROM entreprise._commande) AS commandes_oltp,
  (SELECT COUNT(*) FROM entreprise_dw.faitLivraisons) AS commandes_dw;
\end{verbatim}

\textit{Note :} si le fichier de détails est partiel, on peut avoir plus de commandes que de commandes « avec ventes » (distinct \texttt{nocom} dans \texttt{\_detailcommande}).

\section*{8) Power BI (connexion \& relations)}
\addcontentsline{toc}{section}{8) Power BI (connexion \& relations)}
\textbf{Connexion :} import des tables du schéma \texttt{entreprise\_dw}.

\textbf{Relations (étoile) :}
\begin{itemize}
  \item \texttt{dimTemps[date\_key]} $\rightarrow$ \texttt{faitVentes[date\_key]}
  \item \texttt{dimClient} $\rightarrow$ \texttt{faitVentes} ; \texttt{dimEmploye} $\rightarrow$ \texttt{faitVentes} ;
  \texttt{dimProduit} $\rightarrow$ \texttt{faitVentes} ; \texttt{dimTransporteur} $\rightarrow$ \texttt{faitVentes} ;
  \texttt{dimGeoLivraison} $\rightarrow$ \texttt{faitVentes}
\end{itemize}

\textbf{Dates livraison :} relation active recommandée sur \texttt{faitLivraisons[datecom\_key]}. Les dates \texttt{dateenv\_key} et \texttt{dateobj\_key} peuvent être utilisées via mesures DAX si nécessaire.

\section*{9) Procédure de mise à jour (refresh)}
\addcontentsline{toc}{section}{9) Procédure de mise à jour (refresh)}
\begin{enumerate}
  \item Import/maj CSV dans \texttt{entreprise} (si nouvelles données)
  \item Exécuter le chargement DW (dims puis faits)
  \item Lancer les contrôles (CA, nb commandes)
  \item Power BI : \textbf{Actualiser}
\end{enumerate}

\section*{10) Dépannage rapide}
\addcontentsline{toc}{section}{10) Dépannage rapide}
\begin{itemize}
  \item Erreur FK à l’import : respecter l’ordre parents $\rightarrow$ enfants.
  \item Nouvelles colonnes DW invisibles : vérifier Power Query (\textit{Colonnes conservées}), puis \textit{Fermer \& appliquer}.
  \item Courbes temps peu variées : période courte $\rightarrow$ privilégier analyses produit/client/transporteur.
\end{itemize}

\end{document}